\section{Problem Definition}
\label{sec:problem-definition}

\subsection{Atomic-Core Compilation}

Given $D$, training instances $\mathcal{I}_{\mathrm{train}}$, and normalized
singleton-goal evidence $E$, atomic compilation produces the certified atomic
module core $\mathcal{M}_D$, whose query-independent branches are triggered by
singleton achievements.

\subsection{Query Compilation}

For a validated bound query
$\widehat\tau_q=(\tau_q,\theta_q)$ whose binding maps each parameter to a
type-compatible object, query compilation takes $\mathcal D_q$ and
$\mathcal M_D$ and produces $\mathcal Q_q$: top-level, DFA-dispatch, and
transition-repair plans. A conjunctive guard is compiled into a certified
serialization of signed atomic obligations, not into a trigger in
$\mathcal M_D$. For appended queries $q_1,\ldots,q_k$, the one maintained domain
library is
\[
\mathcal{L}_D^{[0]}=\mathcal{M}_D,
\qquad
\mathcal{L}_D^{[k]}=
\mathcal{M}_D\cup\bigcup_{i=1}^{k}\mathcal{Q}_{q_i}.
\]
Thus $\mathcal M_D$ is the stable core of $\mathcal L_D^{[k]}$, not a
second persisted library, and appending $\mathcal Q_{q_i}$ neither reconstructs
nor reoptimizes the atomic core. Throughout, calligraphic symbols denote the
finite programs, candidate sets, automata, plan sets, and libraries defined
above; $b$ denotes one candidate branch, $e$ one normalized evidence rule, and
$\chi$ one DFA transition guard. The Technical Appendix formalizes the
supported-fragment assumptions and acceptance boundary.
